\documentclass[11pt]{article}
\usepackage[margin=25mm]{geometry}
\usepackage{amsmath,amssymb,amsthm,booktabs,hyperref}
\newtheorem{theorem}{Theorem}
\newtheorem{proposition}{Proposition}
\newtheorem{corollary}{Corollary}
\title{Diagnostic Information Limits under Healthy Drift and Closed-Loop Memory}
\author{Technical working draft; author details to be supplied}
\date{7 October 2026}
\begin{document}
\maketitle
\begin{abstract}
We study information limits for deterministic diagnostic task calendars
in a known stable linear plant with complete retained states and a scalar
healthy input that can drift throughout observation and transition periods.
For fixed Gaussian process noise whose image contains the fault direction,
the constant-force score loss across a missing span determines two regimes:
a sharp leading deficit of order five halves when that loss is positive,
and a deficit of quadratic order when the score is preserved.
The proofs supply calendar-uniform lower bounds and attaining constructions
on the complete input history, while a compact common positive-definite-noise
controller family admits uniform joint asymptotic design.
Controller examples and a separately certified nonlinear mechanical case
identify the information restrictions and additional finite-risk obligations
needed to interpret these limits physically.
\end{abstract}

\paragraph{Draft status.}
This source assembles the independently reviewed positive-definite process-noise
theorem. It is a technical review draft, not an IEEE-formatted submission.
The fixed-controller common-support and quadratic recoverable regime
have separately reviewed proofs and are integrated below.
The declared nonlinear
certificate sweep is still running; only the stated finished cases enter this draft.

\section{Introduction}
Known feedback transforms the fault, the healthy input and the process noise
through the same plant. When every state is retained and the feedback is known,
inverting that common transformation preserves the diagnostic experiment.
Thus a change of output noise correlation alone does not describe the diagnostic
effect of changing a physical controller. Observation blackouts change this
conclusion because the inverse transformation can require an unobserved state.

Task changes also consume the time during which healthy behavior can drift.
We consider two task signs, a known growing fault, and one unrestricted-level
healthy path under each hypothesis. The paths obey a rate limit on the complete
input clock, including every missing slot. The calendar is selected before noise
and has a known terminal horizon. A transition is therefore a missing-data
interval with continued dynamics, rather than a monetary penalty added to a
memoryless observation model.

The quantity of interest is the deficit relative to full records with the healthy
explanation known. Its leading oracle term grows cubically with the horizon.
For positive-definite process noise, the diagnostic deficit grows as the
horizon to the power five halves. A fixed lower-rank noise image can preserve
the constant-force score across a blackout; its optimal deficit then has
quadratic order, even though all latent inputs need not be recoverable. Its exact leading
coefficient contains the actual missing-span information and the healthy drift
rate. The same coefficient is attained by a terminal-balanced calendar with
linearly increasing hold lengths.

Our main contribution is a sharp calendar result for a fixed-dimensional
controlled plant, rather than a new Gaussian nearest-set detector.
The proof handles arbitrary legal calendars in the converse and constructs a
dual statistic in the genuine retained-innovation range in the upper bound.
Two further results explain the role of control: common-plant gains with equal
poles can have different missing-span penalties, and a homogeneous settling
requirement produces a finite set of optimal scalar design candidates.
A fixed bounded task-even bias preserves the positive leading calendar law
and the quadratic order in the recoverable regime.

The information restrictions are part of the problem. The diagnostician reads
only the declared retained states. Additional sensors, the fast state history,
commands that reveal hidden states, or diagnostic readings during the transition
define different experiments. The controller itself has the states needed for
feedback. This distinction isolates a telemetry or admissibility blackout; it
does not claim that feedback stops when diagnostic readings are excluded.

\section{Relation to established results}
Controlled sensing optimizes observation actions for hypothesis testing.
Nitinawarat, Atia and Veeravalli treat known memoryless observation laws and
fixed-sample or sequential tests \cite{controlled}.
Controlled Markovian observations and nonuniform observation costs are also
established \cite{markov}. Their stationary simple-hypothesis and sequential
contracts differ from two complete rate-limited healthy paths, growing input
faults, deterministic horizons and physically missing state intervals.
The present theorem is a second-order information deficit under this declared
contract, not a replacement for those first-order sensing results.

Full-history uncertainty sets and constrained active diagnosis predate this work.
Scott et al. formulate guaranteed diagnosis by separating full output sets under
bounded process and measurement uncertainty \cite{zonotope}.
Their treatment warns against replacing coupled history sets by a product of
marginal sets. Here the healthy path likewise remains one complete path;
Gaussian innovation noise and positive risk tolerances replace bounded-noise
zero-error guarantees.

For a fixed calendar with common Gaussian covariance, closest convex mean sets
and affine detectors supply the statistical interface \cite{convex}.
The input whitening and projection used below are standard linear-Gaussian
tools. The global comparison across calendars, including dense transitions
and unbalanced runs, requires additional arguments because the uncertainty
class is unbounded in its initial level and time advances during missing data.

Symmetric drift compensation and real dead-time optimization also have a long
history. Allan-variance based observing designs explicitly account for elapsed
time and stochastic drift \cite{allan}.
Their drift-spectrum and variance objectives should not be described as lacking
physical dead time. The current deterministic rate adversary and information
objective require a different optimization. These selected comparisons do not
establish exhaustive historical priority; unresolved full-text comparisons remain
in the accompanying literature audit.


Gaussian feedback design is itself established. Kim, Raimondo and Braatz
propagate both means and covariances under affine state feedback and formulate
statistical-distance design with input and state moment constraints \cite{kim}.
Raimondo et al. combine full-history set-valued observers with moving-horizon
input design to retain bounded-uncertainty diagnosis guarantees \cite{closedset}.
These contributions exclude a claim that feedback benefits, common mean/noise
propagation or closed-loop diagnostic design originate here.
Esna Ashari, Nikoukhah and Campbell directly study feedback effects and robust
closed-loop active diagnosis \cite{feedback2012,robustfeedback2012}.
Their publication metadata are verified, but their full theorem coverage remains
an unresolved comparison in this working draft; the current finite audit does
not certify historical exclusion.

Minimax switchback designs with known carryover also optimize switching patterns
under a fixed horizon \cite{switchback}. Their randomized assignment and
worst-case causal-estimation error differ from the deterministic calendars and
two Gaussian drift-path hypotheses here. This is an assumption comparison,
not a claim that switching-pattern minimax theory was absent.

\section{Plant, observation protocol, and information}
Let
\[
x_j=A x_{j-1}+B_u u_j+B(a_j^h+g_j b_j^h)+w_j,\qquad
u_j=-Kx_{j-1},\qquad A_K=A-B_uK .
\]
The dimension is fixed, all matrices and the controller are known,
$B\ne0$, $Q\succ0$, $w_j$ are independent $N(0,Q)$, and $A_K$ is strictly
discrete-time stable. It need not be normal or satisfy $\|A_K\|<1$.
The initial state $x_{-n_0}=0$ is known, $n_0\ge1$ fixed.
Prefix states at $j=-n_0+1,\ldots,0$ are all read at task $g=+1$ and
$a_j=0$. Thereafter the known template is $a_j=\rho_0+\eta j$,
$\rho_0\ge0$, $\eta>0$. Each hypothesis selects its own unrestricted-level
scalar healthy path of rate at most $r/2$, so their actual difference class
is $\mathcal Z_H=\{z\in\mathbb R^{H+n_0}:|\Delta z_j|\le r\}$, $r>0$.
Nature observes the fixed calendar and selects deterministic complete healthy
paths before the independent Gaussian innovations are generated.
The hypotheses are $h\in\{0,1\}$, with $a_j^h=\mathbf1_{\{h=1\}}a_j$.
The planned horizon and calendar are deterministic. Holds read the complete
state at task $g=\pm1$; changes exclude at least $k\ge1$ consecutive state
readings. During excluded slots $g_j$ is known and $|g_j|\le1$.
No state, noise or healthy level is reset.
The diagnostic protocol does not read fast internal states or control commands.
The two task directions admit known exact-$k$ transitions that can be repeatedly
concatenated, and holding and the final retained reads are allowed.
Longer transitions are also admissible in the converse; no terminal task is imposed.
There is no hard actuator, state or tracking budget in this analytical class.
The diagnostic information set omits the internal fast states and the commands
even though the feedback controller can access its state.
One scalar path acts along the fixed physical force vector $B$.

Take a fixed square factor $Q=LL^\top$, $f=L^{-1}B$ and
$F=\|f\|^2=B^\top Q^{-1}B>0$.
Write $a_f=(f a_j)_j$ and $H_gz=(f g_j z_j)_j$.
For a genuine interval between two consecutive retained states, of length $m$,
the conditional innovation has covariance and whitened-input map
\[
W_m=\sum_{r=0}^{m-1}A_K^rQ(A_K^r)^\top\succeq Q,\qquad
H_m=[A_K^{m-1}L,\ldots,L],\qquad T_m=W_m^{-1/2}H_m .
\]
The predictor uses the actual preceding retained state, or known initial state:
\[
x_{t_i}-A_K^m x_{t_{i-1}}
=\sum_{j=t_{i-1}+1}^{t_i}A_K^{t_i-j}B(a_j^h+g_jb_j^h)
+\sum_{j=t_{i-1}+1}^{t_i}A_K^{t_i-j}w_j .
\]
The input intervals use disjoint original innovations; $T_mT_m^\top=I_d$.
Stacking yields $TT^\top=I$, $P=T^\top T$ an orthogonal projector.
The ideal healthy-subtracted complete-record benchmark is $O_H$ below;
it is distinct from the still-composite full-state information. Define
\[
G_{K,H}(\pi)=\tfrac12\inf_{z\in\mathcal Z_H}\|P(a_f-H_gz)\|^2,\qquad
O_H=\tfrac F2\sum_{j=1}^H a_j^2,\qquad
D_{K,H}^*=O_H-\sup_\pi G_{K,H}(\pi).
\]
The full state covariance contains all cross-block terms. The conditional
transformation, rather than an artificial output initialization, gives this
projection. Full-record fault, health and noise share the common invertible
dynamic map, so the whole full-record experiment and its oracle are
control-neutral.

\section{Sharp information deficit}\begin{theorem}[Fixed-order controlled-memory deficit]\label{thm:main}
Fix the preceding contract. For $m=k+1$, define
\[
M_m=\sum_{r=0}^{m-1}A_K^rB,\qquad
\beta_K(k)=mF-M_m^\top W_m^{-1}M_m>0.
\]
Then
\begin{equation}\label{eq:sharp}
D_{K,H}^*=\frac{\sqrt2}{5}\sqrt{F\beta_K(k)r}\,\eta^{3/2}H^{5/2}
+o(H^{5/2}).
\end{equation}
A terminal-balanced symmetric calendar achieves the coefficient at
$\xi_*=2\beta_K(k)\eta/(Fr)$. For every fixed $\xi>0$ it satisfies
\begin{equation}\label{eq:upper}
D_{K,H}(\pi_H^\xi)\le\frac1{10}
\left(\frac{2\beta_K(k)\eta^2}{\sqrt\xi}+Fr\eta\sqrt\xi\right)H^{5/2}
+O(H^2).
\end{equation}
All constants are for the fixed plant, controller and process covariance.
This is leading-coefficient sharpness, not matching $H^2$ terms or a
controller-uniform statement without uniform stability/transient assumptions.
\end{theorem}



\paragraph{Proof mechanism.}
The complete proof is in Appendix~\ref{sec:proof}.
Strict stability and positive-definite process noise make every nonempty
internal-gap coefficient positive. A legal global signed-tent healthy path
then forces a loss on every calendar. The converse pays the same gap loss
only once when balancing the number of runs against total missing time.
For attainment, an observable innovation-range direction is repaired to have
one exact scalar zero-mass constraint. This repair costs only a lower-order
amount even for asymmetric known transition profiles. Genuine adjacent-state
reads at block joins retain all prior state dynamics.

\subsection{Fixed-calendar statistical meaning}
The difference class $\mathcal Z_H$ is a finite-dimensional polyhedron, including
its free initial level. Its linear image in the genuine whitened data coordinates
is a closed polyhedron. The Euclidean closest difference therefore exists.
A minimizing $z^*$ realizes both healthy budgets through
$b^0=z^*/2$, $b^1=-z^*/2$. With $\mathcal B_{r/2}$ the complete healthy class,
the true mean sets are $\mathcal M_0=T H_g\mathcal B_{r/2}$ and
$\mathcal M_1=T a_f+T H_g\mathcal B_{r/2}$.
Choose $\mu_0^*=T H_g(z^*/2)$, $\mu_1^*=T a_f-T H_g(z^*/2)$, and put
$d=\|\mu_1^*-\mu_0^*\|=\sqrt{2G_{K,H}(\pi)}$.
For $d>0$, nearest-pair convexity gives
\[
(\mu_1^*-\mu_0^*)^\top(\mu_0-\mu_0^*)\le0,\qquad
(\mu_1^*-\mu_0^*)^\top(\mu_1-\mu_1^*)\ge0
\]
for every admissible mean under the respective hypothesis.
The statistic
$S=(\mu_1^*-\mu_0^*)^\top(\widetilde Y-\mu_0^*)/d$
has unit variance and mean at most zero under the null, at least $d$ under
the alternative. Thresholding at $z_{1-\alpha}$ has worst false alarm
at most $\alpha$ and worst power at least
$\Phi(d-z_{1-\alpha})$.
The realizing simple pair and the Neyman--Pearson bound show this power is
optimal among tests with the same uniform false-alarm constraint.
Thus, for $0<\alpha,\beta<1$ and $\alpha+\beta<1$, power at least
$1-\beta$ is possible exactly when
\[
G_{K,H}(\pi)\ge\tfrac12(z_{1-\alpha}+z_{1-\beta})^2.
\]
If $d=0$ the closed sets overlap, so no test can have uniform power exceeding
its uniform false-alarm allowance. This is the classical Gaussian closest-set
interface \cite{convex}, not an additional calendar theorem.
The benchmark $O_H$ is the KL information of the ideal complete white record
$(\mathbf1_{\{h=1\}}fa_j+\xi_j)_j$, after ideal subtraction of the healthy input.
It is distinct from the still-composite full-state information, which retains
the healthy-path infimum even when no state is missing.
Write $D_{K,H}(\pi)=O_H-G_{K,H}(\pi)$ for a particular calendar.

% M2S two-regime section fragment v1, 2026-10-07.
% No documentclass; requires amsmath, amssymb, amsthm and the existing
% theorem/proposition environments. All new labels start with m2sreg:.
% Mathematical dependencies, immutable independently accepted sources:
% m2s_singular_memory_theorem_v1.md
% SHA256 8ff2de70274973e4a9e104b26a5bf4be14e0e742998f2c562f6023e6316ab431
% m2s_recoverable_order_v1.md
% SHA256 5b04f87f23fbfe1d2bc1fb1fbb791e783a47e029c7cdf130e1c000dd22d0f4f1
% Their review bindings and finite computational coverage belong in the
% delivery receipt, not in the reader-facing theorem.
\section{Common-noise input images and two deficit regimes}
\label{m2sreg:section}

The positive-definite model admits a fixed-rank extension when the fault
and the healthy input remain in the process-noise image. Consider the
same known initialized plant and diagnostic protocol, with
\begin{equation}\label{m2sreg:plant}
x_j=A_Kx_{j-1}+B\bigl(\mathbf1_{\{h=1\}}a_j+g_jb_j^h\bigr)+G\xi_j,
\qquad \xi_j\overset{\mathrm{iid}}{\sim}N(0,I_p).
\end{equation}
The dimension is fixed, $A_K$ is known and strictly stable,
$G\in\mathbb R^{d\times p}$ has full column rank, and
$0\ne B\in\operatorname{range}G$. Write
$f=G^\dagger B$ and $F=\|f\|^2>0$.
The fixed prefix $j=-n_0+1,\ldots,0$ is fully read, with $g=1,a=0$;
$x_{-n_0}=0$ is known. Afterwards $a_j=\rho_0+\eta j$,
$\rho_0\ge0,\eta>0$. Each hypothesis has its own complete scalar
healthy path, with unrestricted level and rate at most $r/2$, $r>0$.
Thus the actual difference class is
$\mathcal Z_H=\{z:|\Delta z_j|\le r\}$ on every input slot.
The diagnostician reads the retained full states, not the excluded fast
states or control commands. Held tasks have $g=\pm1$ and a reading in
every slot. Moves exclude at least $k\ge1$ slots; their coefficients
$|g_j|\le1$ are known in advance. For the upper bounds, known
exact-$k$ profiles in both directions can be repeatedly concatenated
with holds at any planned start time. Profiles need not be symmetric.
The calendar is chosen before the healthy paths and the innovations.

For each true interval of $m$ input slots between consecutive
retained or initially known states, define
\begin{equation}\label{m2sreg:interval}
H_m=[A_K^{m-1}G,\ldots,G],\quad W_m=H_mH_m^\top,\quad
M_m=\sum_{q=0}^{m-1}A_K^qB,\quad P_m=H_m^\top W_m^\dagger H_m .
\end{equation}
A gap of $k$ excluded states has $m=k+1$ inputs, including its right
retained input. Let $P_\pi$ be the direct sum over the actual input
intervals, with zero on any unobserved terminal inputs, and put
$\theta=(fa_j)_j$, $H_gz=(fg_jz_j)_j$.
Define
\begin{equation}\label{m2sreg:information}
G_{K,H}(\pi)=\frac12\inf_{z\in\mathcal Z_H}
 \|P_\pi(\theta-H_gz)\|^2,\quad
O_H=\frac F2\sum_{j=1}^Ha_j^2,\quad
D_{K,H}^*=O_H-\sup_\pi G_{K,H}(\pi).
\end{equation}
The factor $1/2$ is the KL convention. The supremum does not require
finite-horizon attainment.

\begin{theorem}[Fixed-plant two-regime deficit]\label{m2sreg:theorem}
Under the preceding contract, let
\[
\beta_k=(k+1)F-M_{k+1}^\top W_{k+1}^\dagger M_{k+1}\ge0.
\]
If $\beta_k>0$, then
\begin{equation}\label{m2sreg:positive}
D_{K,H}^*=\frac{\sqrt2}{5}\sqrt{F\beta_k r}\,
 \eta^{3/2}H^{5/2}+o(H^{5/2}).
\end{equation}
The terminal-balanced symmetric calendar attains this coefficient with
$\xi_*=2\beta_k\eta/(Fr)$.
If $\beta_k=0$, there are fixed constants $0<c<C<\infty$ and $H_0$
such that
\begin{equation}\label{m2sreg:zero}
cH^2\le D_{K,H}^*\le CH^2,\qquad H\ge H_0 .
\end{equation}
No sharp $H^2$ constant is asserted. Both statements are for fixed
$A_K,G,B$ and the stated complete scalar path class.
\end{theorem}

\begin{proposition}[Preserved score and finite recovery capacity]
\label{m2sreg:capacity}
Let $\mathcal R_m=\operatorname{span}\{G,A_KG,\ldots,A_K^{m-1}G\}$,
$\mathcal R_\infty=\bigcup_m\mathcal R_m$, and
$\nu_R=\min\{m:\mathcal R_m=\mathcal R_\infty\}$.
Then $\beta_k=0$ if and only if there exists $y\in\mathbb R^d$ with
\begin{equation}\label{m2sreg:score}
G^\top(A_K^\top)^q y=f,\qquad q=0,\ldots,k.
\end{equation}
In that case
\begin{equation}\label{m2sreg:dimension}
k<\nu_R\le\dim\mathcal R_\infty-p+1\le d-p+1,
\qquad k\le d-p.
\end{equation}
This preserves the score of a constant scalar shift in the $f$ direction;
it does not require $P_{k+1}=I_{(k+1)p}$ or recovery of every latent input.
In particular, $p=d$, equivalently $GG^\top\succ0$, forces
$\beta_k>0$ for every $k\ge1$ and recovers the positive-definite theorem.
\end{proposition}

\begin{proof}[Common support, projection, and Proposition~\ref{m2sreg:capacity}]
Let $\mathcal F_{A_K}$ be the initialized causal map from process
inputs to states, and $R$ retain the declared states. Their noise map is
$L_R=R\mathcal F_{A_K}\operatorname{diag}(G)$.
Every mean under either hypothesis is $L_R\vartheta$ with
$\vartheta=(f(a_j^h+g_jb_j^h))_j$. All laws consequently have the
same support $\operatorname{range}L_R$.
For the reduced SVD $L_R=U_s\Sigma_sV_s^\top$, the data on this support
are equivalent to $\Sigma_s^{-1}U_s^\top Y$, with standard Gaussian
noise and mean $V_s^\top\vartheta$. Thus their KL is precisely
$\|V_s^\top(\vartheta_1-\vartheta_0)\|^2/2$.

The invertible block-triangular predictor transform uses the actual
previous retained state:
\[
x_{t_i}-A_K^{t_i-t_{i-1}}x_{t_{i-1}}
=\sum_{j=t_{i-1}+1}^{t_i}A_K^{t_i-j}
 \bigl(B(a_j^h+g_jb_j^h)+G\xi_j\bigr).
\]
It changes $L_R$ into disjoint interval maps $H_m$ and preserves its
rowspace. Therefore $P_m$ is an orthoproject and
\eqref{m2sreg:information} is the exact profiled KL.
This uses rowspace invariance, not an invalid congruence rule for
pseudoinverses. If $U_m$ is an orthonormal basis matrix for $\operatorname{range}W_m$, the
nonredundant white rows are
$(U_m^\top W_mU_m)^{-1/2}U_m^\top H_m$, with identity covariance
of dimension $\operatorname{rank}W_m$.
The full record instead recovers
$G^\dagger(x_j-A_Kx_{j-1})$, giving the stated control-neutral oracle.

We also need a bound valid for all span lengths.
Cayley--Hamilton makes $\mathcal R_d$ invariant and
$\mathcal R_m=\mathcal R_d$ for $m\ge d$.
For such $m$, $W_m\succeq W_d$ with the same kernel, so inverse
monotonicity on $\mathcal R_d$ gives $W_m^\dagger\preceq W_d^\dagger$.
Hence
\begin{equation}\label{m2sreg:gamma}
\Gamma_K=\max_{1\le m\le d}\|W_m^\dagger\|<\infty,\quad
\sup_{m\ge1}\|W_m^\dagger\|\le\Gamma_K,\quad
S_B=\sum_{q\ge0}\|A_K^qB\|<\infty,\quad C_K=\Gamma_KS_B^2 .
\end{equation}
The last summability follows from strict stability, without normality
or $\|A_K\|<1$. The inverse bound does not use $W_1^\dagger$ on
subspaces reached only later.

With $d_m=\mathbf1_m\otimes f$, $H_md_m=M_m$ and
\[
\beta_{m-1}=\|(I-P_m)d_m\|^2.
\]
Zero loss is exactly $d_m\in\operatorname{row}H_m$, which is
\eqref{m2sreg:score}.
The sequence $\beta_\ell$ is nondecreasing: on a common initialized
horizon of $m+1$ constant-shift inputs, retaining times $1,m+1$
has norm loss $\beta_{m-1}$, while deleting the time-$1$ state has
loss $\beta_m$; rowspace inclusion makes loss nondecreasing.
Also $0\le M_m^\top W_m^\dagger M_m\le C_K$, so
\begin{equation}\label{m2sreg:tail}
\beta_\ell\ge(\ell+1)F-C_K,\qquad \beta_\ell/\ell\longrightarrow F.
\end{equation}

Before reachability stabilizes, its dimension grows by at least one:
$\mathcal R_{m+1}=\mathcal R_m$ already implies invariance.
Since $\dim\mathcal R_1=p$, this proves the upper bounds on $\nu_R$.
If \eqref{m2sreg:score} holds, set $\zeta=(A_K^\top-I)y$.
Then $\zeta\perp\mathcal R_k$.
If $k\ge\nu_R$, this annihilates the invariant reached space, and
$G^\top(A_K^\top)^q y$ has zero successive differences for every
$q\ge0$. It remains the nonzero vector $f$, whereas strict stability
makes it tend to zero. This contradiction proves
\eqref{m2sreg:dimension}.
\end{proof}

\begin{proof}[Proof of Theorem~\ref{m2sreg:theorem}]
An unread terminal move can be replaced by continued holding.
The larger experiment dominates the former one, so the supremum can
be restricted to calendars actually reading time $H$.
Their post input spans cover $1,\ldots,H$, with prefix singletons.

\paragraph{Affine spans and the positive-penalty branch.}
For a span of length $m$, write $q_j=j-(m+1)/2$,
$d_q=(q_jf)_j$, $N_m=\sum_{j=1}^m q_jA_K^{m-j}B$.
At center amplitude $A\ge0$ its exact oracle norm loss is
\begin{equation}\label{m2sreg:affine}
\|(I-P_m)(A d_m+\eta d_q)\|^2
=\beta_{m-1}A^2+2\eta A L_{1,m}+\eta^2L_{2,m},
\quad
L_{1,m}=-M_m^\top W_m^\dagger N_m,\quad L_{2,m}\ge0 .
\end{equation}
By \eqref{m2sreg:gamma}, $|L_{1,m}|\le mC_K/2$.
Thus the negative cross terms over any calendar cost at most
\begin{equation}\label{m2sreg:crossbudget}
E_H=C_K\eta(\rho_0+\eta H)H=O(H^2).
\end{equation}
No reflection symmetry is needed.

Here are the precise replacements in the proof of
Theorem~\ref{thm:main}. If $\beta_k>0$,
$b_K=\inf_{\ell\ge k}\beta_\ell/\ell>0$ by monotonicity and
\eqref{m2sreg:tail}. A missing input amplitude is at most twice
its span's center amplitude. Consequently the total oracle norm
loss $T$ obeys
\[
T\ge \frac{b_K}{4}\sum_{j\in\mathcal D_\pi}a_j^2-E_H .
\]
For late macros of length $L=\lfloor H^{3/4}\rfloor$, write
$A$ for the minimum amplitude, $d$ for missing slots, and $M$
for the number of clipped held runs.
The complete signed zero-endpoint tents used in the main proof are
zero on every missing input and its right retained input.
Their scalar shift $\tau=gz$ is therefore zero in every
nonsingleton span, and their full-metric saving is exactly
$F\sum_j(2a_j\tau_j-\tau_j^2)$.
The other bound on the same $T$ is
$\sum_{\rm macros}\beta_kA^2(M-1)-E_H$;
only completely internal gaps are counted here.
Cross-macro missing slots are allocated once in the former bound.

For fixed $\delta\in(0,1)$ split this same global $T$ before grouping
by macros, and put $W=(1-\delta)\beta_k$ and
$C=r(2A-rL/2)/4$. The tent bound, Cauchy and AM--GM give, for $M>0$,
\[
\frac{FC(L-d)^2}{M}+WA^2M+
 \frac{\delta b_K}{4}A^2d-WA^2-2FCL
\ge 2A\sqrt{FCW}(L-d)+
 \frac{\delta b_K}{4}A^2d-WA^2-2FCL .
\]
For $A\ge32FWr/(\delta^2b_K^2)$ the dead-slot term absorbs
the subtracted $d$. When $M=0,d=L$, it gives the same bound directly.
Late amplitudes are of order $H$, so this condition is uniform
in the calendar. Summation leaves
$\sqrt{2FrW}\sum A^{3/2}L$ up to $O(H^{9/4})+E_H$ in twice
the deficit. The Riemann sum over $[\epsilon H,H]$, followed by
$H\to\infty$ and then $\epsilon,\delta\downarrow0$, gives
the lower coefficient in \eqref{m2sreg:positive}.

For completeness, the matching construction also remains observable
at every rank. Its terminal-balanced symmetric blocks have
half-hold $h=n/2$: $h$ positive reads, $k$ missing slots, $n$
negative reads, $k$ missing slots, and $h$ positive reads.
With fixed $\xi>0$, use the main proof's packing in budget $H-1$,
with evenly distributed padding and one to four final positive reads.
Then
\[
M=\sqrt{H/\xi}+O(1),\quad
n_\ell=\xi\ell+O(1),\quad
A_\ell=\eta\xi\ell^2+O(\ell+1).
\]
Each artificial join is a true singleton conditional span using the
previous real state, so no covariance or state initialization is changed.
On an active block $A\ge r(k+n-1)/2$, use
\[
z_i^{+,L}=r(c_0+h-i),\quad
z_i^-=-r[c_0+\min(i-1,n-i)],\quad
z_i^{+,R}=r(c_0+i-1),\qquad c_0=(k+1)/2 .
\]
Set $U_j=f(A-g_jz_j^{\rm loc})$ on holds and $U_j=fA$ on missing
inputs. Inactive blocks, prefix and tail have $U=0$.
The local levels specify an auxiliary dual, not a globally feasible
healthy path. The baseline scalar weights are
$\lambda_j^0=F(g_jA-z_j^{\rm loc})$ on holds and zero on missing inputs.
Their total mass is zero and their exact full-input support is
\[
h_\ell^0=F\left[\frac{Ar}{2}n(n+2k)
 -4r^2\sum_{j=0}^{n/2-1}(c_0+j)^2\right].
\]
Zero missing weights retain the physical partial-sum plateau.

Let $\gamma=(g_jf)_j$, $p=P_\ell\gamma$,
$N=\gamma^\top P_\ell U$, $D=\|p\|^2$ and
$v_\ell=P_\ell U-(N/D)p$.
Full-column $G$ gives $P_1=I_p$. The $2n-2$ held singletons and
contraction on the two fixed-length gaps imply, with $C_k=2(k+2)$,
\[
D\ge F(2n-2)\ge Fn,\quad |N|\le C_kFA,\quad
\gamma^\top v_\ell=0,\quad P_\ell v_\ell=v_\ell,\quad
\|v_\ell-P_\ell U\|^2=N^2/D\le C_k^2FA^2/n .
\]
Indeed each gap contributes at most $(k+1)FA$ in absolute projected
scalar mass, and at most $FA$ in the baseline.
The zero-mass correction
$\lambda_j-\lambda_j^0=\gamma_j^\top(v_\ell)_j-\lambda_j^0$
has $\ell^1$ norm at most $C_k(k+3)FA$.
Its complete path support costs at most
$r(2(n+k)-1)$ times half that norm.
Hence the assembled observable vector has
\[
h_{\mathcal Z}(\lambda)\le\frac{Fr}{2}
 \sum_\ell A_\ell n_\ell^2+O(H^2),\qquad
\|P\theta-v\|^2=O(H^2).
\]
The latter follows from $\|\theta-U\|^2=O(H^2)$ and the summed
repair energy. A span's actual predictor statistic has coefficient
$W_m^\dagger H_mv_m$; its latent coefficient is $P_mv_m=v_m$
and its variance is $\|v_m\|^2$.
Thus the dual uses retained states alone even when $W_m$ is singular.

The two gap input centers are $1/2\pm(n+k)/2$ relative to the
block center. Equation~\eqref{m2sreg:affine} gives their exact
combined norm loss
\[
2\beta_k(A_\ell+\eta/2)^2+
 \frac{\beta_k\eta^2}{2}(n_\ell+k)^2+
4\eta(A_\ell+\eta/2)L_{1,k+1}+2\eta^2L_{2,k+1}.
\]
Its sum is $2\beta_k\sum A_\ell^2+O(H^{3/2})$.
The full-input dual inequality consequently gives
\begin{equation}\label{m2sreg:symmetricupper}
D_{K,H}(\pi_H^\xi)\le
\frac1{10}\left(\frac{2\beta_k\eta^2}{\sqrt\xi}
 +Fr\eta\sqrt\xi\right)H^{5/2}+O_\xi(H^2).
\end{equation}
Optimization in fixed $\xi$ proves \eqref{m2sreg:positive}.
All rank-sensitive replacements have now been supplied: common
support, the real projection, fixed-plant inverse bounds, and
singleton coercivity for the scalar mass repair.

\paragraph{Zero penalty: a profile-aware upper bound.}
Assume $\beta_k=0$ and fix $L>n_0+k+2$.
Keep the full-input cumulative task coefficient
$M_t=\sum_{j=-n_0+1}^t g_j$.
Hold the current sign $s$ until $sM_t\ge L$.
If at least $k+2$ slots remain, execute the known exact-$k$ move,
flip $s$, and continue; otherwise hold to $H$.
This is a deterministic calendar using the known profile at each
planned start. It balances all input coefficients, including the
missing ones, rather than just held signs.

The threshold overshoot is less than one and a move changes $M_t$
by at most $k$. Thus the next opposite hold run reaches its threshold
after between $2L-k$ and $2L+k+2$ slots.
The initial run exceeds two slots; a skipped final move adds at most
$k+1$ holds. For large $H$, every post hold run has at least two
slots and
\begin{equation}\label{m2sreg:balance}
|M_t|\le L+k+2 .
\end{equation}
If $N_h,N_m$ count post held inputs and moves, then
$H=N_h+kN_m$ and $N_h\ge2N_m$.
The singleton count, excluding the right retained point of each gap,
is $S=N_h-N_m\ge H/(k+2)$.

Put $m=k+1$. Since $P_md_m=d_m$, only the affine slope is lost
in each fixed-length gap:
\begin{equation}\label{m2sreg:slopeloss}
E=\|(I-P)\theta\|^2\le
N_m\eta^2F\,m(m^2-1)/12=O(H).
\end{equation}
Let $\gamma=(g_jf)_j$, $w=P\theta$, $p=P\gamma$ and
$\lambda_j=\gamma_j^\top w_j$.
We have $\lambda_j=Fg_ja_j+e_j$, where $e_j=0$ on singletons
and is uniformly bounded on the fixed-length gaps independently
of the amplitude. The post cumulative $S_t=\sum_{j=1}^tg_j=M_t-n_0$
has $|S_t|\le L+k+2+n_0$ by \eqref{m2sreg:balance}. Abel's identity gives
\[
\sum_{j=1}^t g_ja_j=a_tS_t-\eta\sum_{j=1}^{t-1}S_j .
\]
All partial $\lambda$ masses are therefore $O(t+1)$ and
$N=\sum_j\lambda_j=O(H)$. The fixed prefix changes only the constants.

Now $D=\|p\|^2$ lies between $FS\ge FH/(k+2)$ and $F(H+n_0)$.
Set $\alpha=N/D=O(1)$ and $v=w-\alpha p$.
Then $Pv=v$, $\sum_j\gamma_j^\top v_j=0$, and
\begin{equation}\label{m2sreg:globalrepair}
\|\theta-v\|^2=E+\alpha^2D=O(H).
\end{equation}
Each $\gamma_j^\top p_j$ is bounded by $F\sqrt m$ because
all nonsingleton spans have fixed length. The repaired scalar weights
$\lambda'_j=\gamma_j^\top v_j$ hence still have partial masses $O(t+1)$.
The free-level complete Lipschitz class has the exact support formula
\begin{equation}\label{m2sreg:support}
h_{\mathcal Z}(\lambda')=
r\sum_{t=-n_0+1}^{H-1}
 \left|\sum_{j=-n_0+1}^t\lambda'_j\right|=O(H^2).
\end{equation}
Summation by parts proves this formula; a nonzero total mass instead
has infinite support because the initial level is free.
The observable range vector $v$ yields
\begin{equation}\label{m2sreg:dual}
G_{K,H}(\pi)\ge
 \langle v,\theta\rangle-\tfrac12\|v\|^2-h_{\mathcal Z}(\lambda'),
\qquad
O_H-G_{K,H}(\pi)\le
 \tfrac12\|\theta-v\|^2+h_{\mathcal Z}(\lambda').
\end{equation}
Equations~\eqref{m2sreg:globalrepair}--\eqref{m2sreg:dual}
prove the $O(H^2)$ upper bound.

\paragraph{Zero penalty: an all-calendar lower bound.}
By monotonicity and \eqref{m2sreg:tail}, define finite
\[
L_z=\max\{m\ge1:\beta_{m-1}=0\}\ge k+1,\qquad
b_+=\inf_{m>L_z}\frac{\beta_{m-1}}m>0 .
\]
The capacity bound gives $L_z\le\nu_R$.
Set $\epsilon=(4L_z+2)^{-1}$ and $J_0=\lceil\epsilon H\rceil$.
Consider any calendar reading $H$.

If a positive-penalty interval ends at $t\ge J_0$, it lies entirely
in the post affine region and has center amplitude
$A\ge\eta(t+1)/2\ge\eta\epsilon H/2$, even if its span is long.
For sufficiently large $H$, uniformly in the span,
\eqref{m2sreg:affine} gives
$\|(I-P_m)\theta_m\|^2\ge(b_+/2)mA^2$.
Using $z=0$ in the profile therefore gives
\begin{equation}\label{m2sreg:positivecase}
O_H-G_{K,H}(\pi)\ge\tfrac12\|(I-P)\theta\|^2
\ge \frac{b_+\eta^2\epsilon^2}{16}H^2 .
\end{equation}

Otherwise every interval ending at or after $J_0$ has length
at most $L_z$ and preserves $d_m$.
At least $(1-\epsilon)H/L_z-O(1)$ such late intervals cover
inputs $J_0,\ldots,H$. Their center amplitudes are at least
$\eta\epsilon H/2$ for large $H$.
Choose the single complete healthy difference
\[
z_j=(r/2)g_j,\qquad \widetilde h_j=fg_jz_j=(r/2)fg_j^2 .
\]
Because $|g_j-g_{j-1}|\le2$, this path obeys $|\Delta z|\le r$,
and the two paths $z/2,-z/2$ meet their individual rate bounds.
The free prefix level makes the same choice legal there.
For a late zero-penalty span, preservation of $d_m$ gives
\[
\langle P_m\theta_m,P_m\widetilde h_m\rangle
=\frac r2FA\sum_{j=1}^m g_j^2+
 \eta\langle P_md_q,P_m\widetilde h_m\rangle
\ge \frac r2FA-C_0 .
\]
Here $C_0$ is fixed since $m\le L_z$.
The right retained input has $g^2=1$, so $\sum g_j^2\ge1$.
No entrywise positivity of the projection is assumed.
The total late cross is thus at least
\[
\frac{rF\eta\epsilon(1-\epsilon)}{4L_z}H^2-O(H).
\]
All early intervals end before $J_0$.
Contraction and Cauchy bound their cross from below by
$-(r/2)F\sum_{\rm early}\sqrt{m\sum a_j^2}$.
Their total input count is at most $J_0-1+n_0$, and
$a_j\le\rho_0+\eta J_0$, so the early contribution is at least
$-(r/2)F\eta\epsilon^2H^2-O(H)$.
Finally $\|P\widetilde h\|^2\le r^2F(H+n_0)/4=O(H)$.
Evaluating the profiled infimum at this legal full path and expanding
the half square gives
\begin{align}
O_H-G_{K,H}(\pi)
&\ge \tfrac12\|(I-P)\theta\|^2+
 \langle P\theta,P\widetilde h\rangle-
 \tfrac12\|P\widetilde h\|^2 \notag\\
&\ge \frac{rF\eta}{2}
 \left[\frac{\epsilon(1-\epsilon)}{2L_z}-\epsilon^2\right]H^2-O(H)
 =\frac{rF\eta\epsilon}{8L_z}H^2-O(H).
\label{m2sreg:zerocase}
\end{align}
The last equality uses $\epsilon=(4L_z+2)^{-1}$.
All constants in both cases are independent of the calendar.
A smaller positive constant than the minimum of the coefficients
in \eqref{m2sreg:positivecase} and \eqref{m2sreg:zerocase}
gives the lower bound in \eqref{m2sreg:zero}.
Together with the profile-aware upper construction, this completes
the two-regime theorem.
\end{proof}

The distinction concerns recoverability of the constant force score.
A zero coefficient does not remove healthy uncertainty: with
$r,\eta>0$ and held $g=\pm1$, its optimal deficit still has order $H^2$.
Positive-definite noise excludes the zero-score-loss case, whereas
a fixed lower-rank noise image can preserve that score across a gap.
Both conclusions use the same physical input transformation for the
fault, the complete healthy path, and the noise.
% End of the fixed-plant section. No controller-family uniformity claim,
% partial-observation model, or additional theoretical direction is added.


\section{Control consequences}
\subsection{Equal poles do not fix diagnostic loss}
Take the same physical plant and process noise
\[
A=\operatorname{diag}(4/5,3/5),\quad B_u=I_2,\quad B=(1,2)^\top,
\quad Q=\begin{pmatrix}1&1/3\\1/3&5/9\end{pmatrix}.
\]
Here $F=29/4$. Gains
$K_d=\operatorname{diag}(3/10,7/20)$ and
$K_n=\left(\begin{smallmatrix}3/10&-2\\0&7/20\end{smallmatrix}\right)$
give closed matrices $\operatorname{diag}(1/2,1/4)$ and
$\left(\begin{smallmatrix}1/2&2\\0&1/4\end{smallmatrix}\right)$.
They share the poles and full oracle, but at $k=1$,
\[
\beta_{K_d}(1)=1343/344,\qquad \beta_{K_n}(1)=18007/10456.
\]
The ratio of sharp deficit coefficients is approximately $0.664$.
This is an ideal known-gain comparison with common $B,Q$ and no actuator
budget; it does not identify a hardware-optimal controller.

\subsection{A homogeneous settling constraint}
For a scalar common-input plant with $Q=q>0$, $B=1$ and closed parameter
$0\le c<1$, write
\[
A_m(c)=\sum_{j=0}^{m-1}c^j,\quad B_m(c)=\sum_{j=0}^{m-1}c^{2j},
\quad \kappa_c(k)=k+1-A_{k+1}(c)^2/B_{k+1}(c).
\]
Thus $F=1/q$, $\beta=\kappa/q$ and the coefficient is
$\sqrt{2}\sqrt{\kappa r}\eta^{3/2}/(5q)$.
At fixed $k\ge1$, $\kappa$ strictly decreases in $c\in[0,1)$.
To impose a common homogeneous qualification, let $|e_0|\le E$,
$e_{j+1}=ce_j$, $\gamma=\epsilon/E\in(0,1)$ and a fixed integer
move duration $d_{\rm move}\ge0$. Require the smallest positive $s$ with
$c^s\le\gamma$, and set $k(c)=d_{\rm move}+s(c)$.
This requirement concerns the homogeneous component; forcing and innovations
continue during the excluded slots. It does not assert full noisy return to a tube.

\begin{proposition}[Plateau candidates]\label{prop:settling}
On a compact $0<c_{\rm lo}\le c\le c_{\rm hi}<1$, the minimum of
$\kappa_c(k(c))$ is achieved in the finite set
\[
\{c_{\rm hi}\}\ \cup\
\{\gamma^{1/s}:c_{\rm lo}\le\gamma^{1/s}\le c_{\rm hi},
\ 1\le s\le s(c_{\rm hi})\}.
\]
If $L=-\log\gamma$ and $\delta=-\log c$, then
\[
\delta\kappa_c(k(c))\longrightarrow L-2\tanh(L/2)>0
\quad(c\uparrow1).
\]
\end{proposition}
\begin{proof}
For fixed $m=k+1$ the explained constant-mode norm is
$R_m=(1+c)(1-c^m)/[(1-c)(1+c^m)]$.
Its logarithmic derivative is positive because
$\sum_{j=0}^{m-1}c^{2j}>mc^{m-1}$ by strict AM--GM.
The exact increment in $k$ is
\[
\kappa_c(k+1)-\kappa_c(k)
=\frac{(B_m-A_mc^m)^2}{B_m(B_m+c^{2m})}>0.
\]
The plateau $s\ge2$ is $(\gamma^{1/(s-1)},\gamma^{1/s}]$;
the first plateau has right endpoint $\gamma$.
Each nonempty intersection with the compact controller interval contains its
rightmost point. Fixed-$k$ monotonicity places its minimum there.
At the equality $c^s=\gamma$ the shorter $s$ applies, so a floating-point
logarithm must not change the tie rule.
Finally $m\delta=L+O(\delta)$ and
$R_m=\coth(\delta/2)\tanh(m\delta/2)$.
The limit follows, and positivity is
$L-2\tanh(L/2)=\int_0^L\tanh^2(t/2)\,dt>0$.
\end{proof}
For $\gamma=81/100$ and $d_{\rm move}=0$, the finite controller family
$\{4/5,9/10,19/20\}$ has $k=1,2,5$ and exact gap penalties
$1/41$, $2/91$, $5064645/111045881$.
The middle gain is best in that family. On the full interval $[4/5,19/20]$,
the unique best is instead $c=81/100$, with $\kappa=361/16561$.
The latter ranking uses rational isolation of the irrational plateau endpoints.
It optimizes the declared single-gap or coefficient function. Pointwise
calendar sharpness alone does not exchange the limit with optimization over
an entire continuous controller family. For a finite fixed-gain family,
the maximum of finitely many little-oh remainders is still little-oh.


\subsection{A uniform bridge to joint asymptotic design}
\begin{theorem}[Compact common-noise controller family]\label{thm:uniform}
Let the fixed plant, $B\ne0$, $Q\succ0$, healthy rate, fault template,
prefix and diagnostic information be common. Let $\mathcal K$ be a nonempty
compact gain set for which each $A_K$ is strictly Schur, and let
$k(K)\in\{1,\ldots,k_{\max}\}$. For each gain assume repeatedly concatenable
known exact-$k(K)$ transitions in both directions, bounded by $|g|\le1$.
Set $C(K)=\sqrt2\sqrt{F\beta_K(k(K))r}\eta^{3/2}/5$. Then
\[
\sup_{K\in\mathcal K}|H^{-5/2}D_H^{*,K}-C(K)|\longrightarrow0,
\]
and the joint deficit divided by $H^{5/2}$ tends to $\inf_K C(K)$.
\end{theorem}
\begin{proof}
Continuity of the spectral radius on the compact family gives
$a_*:=\max_K\rho(A_K)<1$. Choose $a_*<\tau<1$.
The resolvent has a finite norm maximum $R_*$ on
$\mathcal K\times\{|z|=\tau\}$. Matrix Cauchy integration gives
$\|A_K^j\|\le R_*\tau^{j+1}$, hence a common finite bound
$S_*\ge\sum_j\|A_K^jB\|$.
Since $W_{m,K}\succeq Q$, the affine-cross constant in the main proof is
uniformly bounded by $C_*:=\|Q^{-1}\|S_*^2$.

Each $\beta_K(\ell)$ is continuous and strictly positive for fixed
$1\le\ell\le k_{\max}$. Finite $k$ and compactness give common
$0<\beta_{\rm lo}\le\beta_K(k(K))\le\beta_{\rm hi}$.
The tail $\beta_K(\ell)\ge(\ell+1)F-C_*$ and monotonicity give
a common positive $b_*$ below all $\beta_K(\ell)/\ell$, $\ell\ge1$.
For fixed $\varepsilon,\theta$, the converse proof therefore uses one
absorption threshold and one $o(H^{5/2})$ remainder for all gains and calendars.
Its lower coefficient is
$\sqrt{1-\theta}(1-\varepsilon^{5/2})C(K)$.
First choose $\varepsilon,\theta$ using $\sup C(K)<\infty$, then the common
horizon threshold; this gives the uniform lower comparison.

The optimal $\xi_K=2\beta_K(k(K))\eta/(Fr)$ lies in a fixed positive compact
interval. Packing has a common bounded padding, because the unused time is
less than the next block length and is evenly distributed.
Thus $n_l=\xi_Kl+O(1)$ and $A_l=\eta\xi_Kl^2+O(l+1)$ uniformly.
A common finite early cutoff makes every later block active: a lower quadratic
$\eta\xi_{\rm lo}l(l-1)$ eventually dominates the active threshold's common
upper linear bound. Auxiliary error, repair energy and support correction
are all bounded by common multiples of
$\sum(n_l+k)^3$, $\sum A_l^2/n_l$ and $\sum A_l(n_l+k)$, each $O(H^2)$.
The fixed spans have common affine coefficients, and the two power sums in
the main proof have common $O(H^2)$ errors. Hence
$D_H^{*,K}\le C(K)H^{5/2}+C_UH^2$ for one $C_U$.
This proves uniform convergence. Taking infima changes any uniformly bounded
error by no more than its uniform bound, yielding the joint statement.
\end{proof}
No continuity of $k(K)$ or of its known profiles is needed for this limit;
neither ensures a finite-horizon optimizer or a minimum of $C(K)$.
In the declared scalar settling example, the finitely many plateau endpoints
are lower semicontinuous because the shorter $k$ is included at each equality.
Its unique coefficient minimum is $c_*=81/100$.
Theorem~\ref{thm:uniform} therefore makes that fixed controller and its
terminal-balanced calendar jointly optimal at leading order under the same
homogeneous qualification. Any sequence of controllers within
$o(H^{5/2})$ of the optimal deficit converges to $c_*$: on the compact set
outside each neighborhood, lower semicontinuity and uniqueness give a positive
coefficient gap, while the information remainder is uniform.
This adds a joint asymptotic conclusion, not a computed finite-$H$ risk optimum
or complete forced-return guarantee.

\section{Bounded task-even uncertainty}
\begin{corollary}\label{cor:even}
Add $Be_j^h$ with $|e_j^h|\le E/2$ to the input, where $E$ is fixed.
Assume the complete difference class is the explicit Minkowski envelope
$g\mathcal Z_H+\mathcal E_H$, with $0\in\mathcal E_H$ and
$|e_j|\le E$, and leave the covariance, plant and calendar class unchanged.
With $N=H+n_0$ and $A_H=(\sum a_j^2)^{1/2}$, every calendar satisfies
\[
0\le G_{\rm odd}(\pi)-G_{\rm mix}(\pi)
\le FE\sqrt N A_H=O(H^2).
\]
The same bound holds for $D_{\rm mix}^*-D_{\rm odd}^*$ under the same oracle.
Thus the leading coefficient of Theorem~\ref{thm:main} is unchanged.
\end{corollary}
\begin{proof}
In the real input projection let $t=P(fa)$ and
$\mathcal K=\{P(fgz):z\in\mathcal Z_H\}$.
Each extra mean $P(fe)$ has norm at most $R=E\sqrt{FN}$.
For $d=\operatorname{dist}(t,\mathcal K)$, class inclusion and the triangle
inequality give $(d-R)_+\le d_{\rm mix}\le d$.
Since $0\in\mathcal K$, $d\le\sqrt F A_H$.
For $G=d^2/2$, separately considering $d\ge R$ and $d<R$ gives
$0\le G_{\rm odd}-G_{\rm mix}\le dR\le FE\sqrt N A_H$.
Taking suprema over the same calendars proves the deficit comparison.
\end{proof}
This is an order-two perturbation bound, not an exact constant translation.
Multiplicative uncertainty that changes the plant, covariance or support is a
different model. If an unrestricted even difference class contains $a$,
choosing $e=a$ and $z=0$ makes the hypotheses indistinguishable.

\section{Verification and a separate nonlinear comparator instance}
Exact rational checks compare the initialized state covariance against the
retained-input projector in scalar and matrix examples. The matrix report
contains 150 retained-subset checks and 68 full calendars, including rotating,
nonnormal and Jordan models. The proof is universal; these finite checks
support implementation. The symmetric transition profiles in the original
formal reports have zero total repair, so they do not numerically exercise
nonzero $\alpha$. A separately reviewed asymmetric matrix example has
$N=-22863/1307$, $D=628321/10456$ and $\alpha=-182904/628321$.
Its repaired healthy mass is exactly zero. The separately reproduced
six-case asymmetric supplement includes two low-rank examples with repair
energies $100/77$ and $507/178$. These supplement the original zero-repair
checks without changing their historical counts.

The separate planar six-joint mechanical certificate uses a fixed joint-4
fault $f(t)=0.0003(t-1)_+$, four healthy prefix slots, fast step $0.001$ s,
and diagnostic slot $0.25$ s. Move and settle exclude twelve slots while
all process and encoder innovations continue. Point readout is the last
existing pre-step fast sample in the slot; average readout uses its 250 samples.
The declared families and rounding were frozen before the scan.
This comparator has parameter-dependent physical bounds and deterministic
nonlinear approximation errors; it is not claimed to be a direct exact
instance of the fixed common-covariance matrix theorem.

For the finished terminal-balanced $H=120$ case, complete-history variance
certificates and two-sided mean/error/event budgets give
\begin{center}
\begin{tabular}{lccc}
\toprule
Readout & Upper variance & Protected gap & Power lower bound\\
\midrule
250-sample average & $6.7073942\,10^{-7}$ & $0.0111325$ & $0.999999999979$\\
Last existing point & $0.0181313$ & $-0.453312$ & $0$\\
\bottomrule
\end{tabular}
\end{center}
Both false-alarm upper bounds are approximately $0.000967122592$.
The point row fails the sufficient power requirement; its lower bound zero is
not a statement that actual power is zero. The inherited equilibrium enclosure
has been independently accepted for the frozen ramp through $H\le600$
and actual $f\le0.045$, not for every signal in the wider historical force box.
The 400-row sweep is incomplete, and no all-family shortest-time or sequential
optimality conclusion is asserted here. Full certificates retain failed,
degenerate and not-run rows with their identities.

\section{Limits and conclusions}
The calendar laws assume a known fixed stable plant and controller, complete
retained states, a fixed full-column Gaussian noise factor and a nonzero scalar
fault/healthy force in its image. The compact controller-family theorem
additionally retains common positive-definite process noise.
It does not cover noisy partial measurements, unknown matrices, arbitrary
vector nuisance, saturation, adaptive task calendars or unknown fault onset.
Near-singular or rank-changing controller families require separate uniform
analysis. Finite healthy amplitude boxes have a different long-horizon regime.
The common-support two-regime extension is a fixed-controller result.
The complete finite-risk sweep, hardware model adaptation and final manuscript
acceptance remain unfinished in the current research workflow.

The result identifies a control-dependent missing-span penalty within a
common physical experiment and proves how it interacts with full-history
healthy drift, including the distinction between losing and preserving the
constant-force score. The matching calendar gives an explicit design rule; the
controller examples explain which parts of that rule survive common input
and timing constraints. The separate mechanical certificate illustrates the
additional obligations required for finite physical risk statements.

\appendix
\section{Complete proof of Theorem~\ref{thm:main}}\label{sec:proof}

\begin{proof}
\paragraph{Strict gap loss and its tail.}
For one input span of length $m$, set $P_m=H_m^\top W_m^{-1}H_m$ and
$d_m=\mathbf1_m\otimes f$.
Its squared-input-information loss at constant amplitude $a$ is
$a^2d_m^\top(I-P_m)d_m=\beta_K(m-1)a^2$.
If $m\ge2$ and this coefficient were zero, then $d_m\in\operatorname{range}H_m^\top$.
Thus $H_m^\top y=d_m$ for some $y$. Its last two input blocks require
$L^\top y=f$ and $L^\top A_K^\top y=f$.
The first gives $y=Q^{-1}B\ne0$; the second gives
$A_K^\top Q^{-1}B=Q^{-1}B$, an eigenvalue-one contradiction.
Hence $\beta_K(\ell)>0$ for $\ell\ge1$.
Nested state deletion in the same initialized full experiment shows
$\beta_K(\ell)$ is nondecreasing in $\ell$.

Strict stability gives
$S_B=\sum_{r\ge0}\|A_K^rB\|<\infty$.
Set $C_K=\|Q^{-1}\|S_B^2$.
Since $\|M_m\|\le S_B$ and $W_m^{-1}\preceq Q^{-1}$,
$\beta_K(\ell)\ge(\ell+1)F-C_K$ and $\beta_K(\ell)/\ell\to F$.
Thus $b_K=\inf_{\ell\ge k}\beta_K(\ell)/\ell>0$.
For example, with $L_0=\max(k,\lceil2C_K/F\rceil)$,
monotonicity and the tail bound give
$b_K\ge\min(F/2,\beta_K(k)/L_0)>0$.

\paragraph{Affine non-reflection cross.}
Let $h_j=j-(m+1)/2$, $d_h=(h_jf)_j$, and
$N_m=H_m d_h=\sum_jh_jA_K^{m-j}B$.
The centered constant--slope loss coefficient is
\[
L_{1,m}=d_m^\top(I-P_m)d_h=-M_m^\top W_m^{-1}N_m .
\]
The unprojected inner product is zero because $\sum_jh_j=0$.
Also $L_{2,m}=d_h^\top(I-P_m)d_h\ge0$ and
$\|N_m\|\le(m/2)S_B$, so $|L_{1,m}|\le mC_K/2$.
For affine input $x_j=(A+\eta h_j)f$, $A\ge0$,
\begin{equation}\label{eq:affine}
\| (I-P_m)x\|^2
=\beta_K(m-1)A^2+2\eta A L_{1,m}+\eta^2L_{2,m}
\ge\beta_K(m-1)A^2-C_K\eta A m .
\end{equation}
No reflection or normality has been assumed.
The true non-singleton intervals are disjoint; their total length is at most
$H$ and their center amplitudes at most $\rho_0+\eta H$.
Consequently their total negative cross cost is bounded uniformly over
calendars by $E_H=C_K\eta(\rho_0+\eta H)H=O(H^2)$.

\paragraph{Full-history calendar-uniform converse.}
Replacing terminal moves without later readings by holding preserves the old
data and adds readings. Data processing for each full healthy pair and then
an infimum shows that the optimal calendar can end with a retained state
at $H$. Prefix states are all read, so every non-singleton interval is
post-onset and internal.

Fix $\varepsilon\in(0,1)$ and divide the late interval into complete macros
of length $L=\lfloor H^{3/4}\rfloor$.
On each clipped observed run of length $n$ use one scalar signed tent
$z_i=g_i r\min(i,n-1-i)$.
Set it to zero on missing inputs, right-first retained inputs, run endpoints,
prefix, macro joins and unused tail. This is one full rate-$r$ path, realized
by $z/2,-z/2$ under the two hypotheses.
Writing $h=gz$, every non-singleton input span has residual $a-h=a$.
If $T$ is the sum of its actual gap norm losses, orthogonal decomposition gives
\begin{equation}\label{eq:converse}
2D_{K,H}(\pi)\ge F\sum_j(2a_jh_j-h_j^2)+T .
\end{equation}
No stationary endpoint or prefix replacement is needed.

In a macro let its minimum amplitude be $A$, its number of missing slots $d$,
its clipped run count $M$ and lengths $n_i$, with $\sum n_i=L-d$.
Put $C=r(2A-rL/2)/4>0$ for large $H$.
Tent area and height give saving at least $C\sum n_i^2-2CL$.
There are $M-1$ complete actual gaps wholly within the macro, each with length
at least $k$ and center amplitude at least $A$.
Equation \eqref{eq:affine} and monotonicity give
\[
T\ge\sum_{\rm macros}\beta_K(k)A^2(M-1)-E_H.
\]
For each gap its missing amplitudes are at most twice its input-center
amplitude; hence the same $T$ also satisfies
\[
T\ge(b_K/4)\sum_{\rm macros}A^2d-E_H .
\]
Crossing gaps contribute no complete-gap count but each missing node is
assigned once in this second budget.

Fix $\theta\in(0,1)$ and split the same global $T$ before applying its bounds.
Set $W=(1-\theta)\beta_K(k)$.
For $M>0$, Cauchy and AM--GM lower-bound a macro contribution by
\[
FC(L-d)^2/M+WA^2M+\theta b_KA^2d/4-WA^2-2FCL
\ge2A\sqrt{FCW}(L-d)+\theta b_KA^2d/4-WA^2-2FCL.
\]
The condition $A\ge32FWr/(\theta^2b_K^2)$ absorbs its missing-slot term,
uniformly on late macros. For $M=0,d=L$, the dead budget gives the same
bound $2A\sqrt{FCW}L-WA^2-2FCL$ without division by zero.
The cross error is paid once after the convex split.

As $A=\Theta(H)$ and $C=(rA/2)(1+O(H^{-1/4}))$, the leading sum is
$\sqrt{2FrW}\sum A^{3/2}L$.
Coefficient, rounding and incomplete Riemann-block errors are $O(H^{9/4})$,
$\sum A^2=O(H^{9/4})$, $\sum CL=O(H^2)$ and $E_H=O(H^2)$.
All constants are calendar-independent.
The Riemann sum equals
$(2/5)\eta^{3/2}(1-\varepsilon^{5/2})H^{5/2}+o(H^{5/2})$.
Taking the optimal deficit before the limit and then
$\varepsilon,\theta\downarrow0$ proves
\begin{equation}\label{eq:lower}
\liminf D_{K,H}^*/H^{5/2}\ge
\frac{\sqrt2}{5}\sqrt{F\beta_K(k)r}\,\eta^{3/2}.
\end{equation}

\paragraph{Packing and genuine singleton joins.}
A symmetric block with even $n=2h$ has $h$ reads at $+$, $k$ missing slots,
$n$ reads at $-$, $k$ missing slots and $h$ reads at $+$.
With budget $H-1$, fit the maximal number of blocks with
$\bar n_\ell=2\lceil\xi\ell/2\rceil$; evenly distribute remaining multiples of
four slots, adding two to $n_\ell$ per unit and using a fixed left-to-right
tie rule. Observe the last one to four states at $+$.
Every block padding is $O(1)$ and
\[
M=\sqrt{H/\xi}+O(1),\quad n_\ell=\xi\ell+O(1),\quad
A_\ell=\rho_0+\eta(x_\ell+c_\ell)=\eta\xi\ell^2+O(\ell+1).
\]
At an artificial join both adjacent states are actually read, so the next
conditional input interval is a singleton.
The global input projector therefore splits into its genuine block spans.
It does not reset any state or make output covariances independent.

\paragraph{Vector auxiliary and one scalar weighted constraint.}
On active blocks, $A\ge r(k+n-1)/2$, use $h=n/2$, $c_0=(k+1)/2$ and local levels
\[
z_i^{+,L}=r(c_0+h-i),\quad
z_i^-=-r[c_0+\min(i-1,n-i)],\quad
z_i^{+,R}=r(c_0+i-1).
\]
Held white auxiliary inputs are $U_j=f(A-g_jz_j^{loc})$ and missing ones $fA$.
On inactive blocks, prefix and tail set $U=0$.
Only uniformly finitely many early blocks are inactive.
Local levels generate a dual rather than a claimed global primal path.

The baseline scalar direction has $\lambda_j^0=(g_jf)^\top U_j$ on holds
and zero on missing inputs. It has zero mass and support
\begin{equation}\label{eq:baseline}
h_\ell^0=F\left\{\frac{Ar}{2}n(n+2k)
-4r^2\sum_{j=0}^{n/2-1}(c_0+j)^2\right\}.
\end{equation}
Partial sums have the correct left-positive/right-negative signs,
are zero at the negative center and at the end, and the local levels
saturate $\Delta z=-r\Delta t\,\operatorname{sign}(Q)$ wherever $Q\ne0$.
Missing zero weights preserve the physical gap span $k+1$.
Summation by parts and fourfold layer counting prove \eqref{eq:baseline}.

Write $h_\ell=(g_jf)_j$, $p_\ell=P_\ell h_\ell$, and
\[
N_\ell=h_\ell^\top P_\ell U_\ell,\quad
D_\ell=\|p_\ell\|^2,\quad \alpha_\ell=N_\ell/D_\ell,\quad
v_\ell=P_\ell U_\ell-\alpha_\ell p_\ell .
\]
Apart from two gap-right inputs, the other $2n-2$ held inputs are singletons
with projection $I_d$ and $\|h_j\|^2=F$.
Thus $D_\ell\ge F(2n-2)\ge Fn$.
The difference $(h_j^\top(P_\ell U)_j)_j-\lambda^0$ lies in two fixed
spans only. In each span,
$\sum_j|h_j^\top(PU)_j|\le\sqrt{mF}\|PU\|\le mFA$;
the baseline right-held point contributes at most $FA$.
With $C_k=2(k+2)$ this gives
\[
|N_\ell|\le C_kFA,\quad |\alpha_\ell|\le C_kA/n,\quad
h_\ell^\top v_\ell=0,\quad P_\ell v_\ell=v_\ell,\quad
\|v_\ell-P_\ell U_\ell\|^2=N_\ell^2/D_\ell\le C_k^2FA^2/n .
\]
The repaired mass is one scalar constraint, exactly matching one unrestricted
scalar healthy level. It is not a coordinatewise vector-nuisance constraint.

For $\delta\lambda_j=h_j^\top v_j-\lambda_j^0$,
$\sum\delta\lambda_j=0$ and
$\sum|h_j^\top p_j|\le\|h_\ell\|\|p_\ell\|\le L_\ell F$.
Hence $\|\delta\lambda_\ell\|_1\le C_k(k+3)FA$.
The full scalar rate class has support $r\sum|Q_j|$ for a zero-mass vector;
each partial sum is at most half its $L_1$ norm.
A length-$L_\ell$ block correction therefore costs
at most $r(L_\ell-1)\|\delta\lambda_\ell\|_1/2=O(FrA(n+k))$.
Summation and restriction of the full path class give
\begin{equation}\label{eq:support}
h_{\mathcal Z}(\lambda)\le\sum h_\ell^0+O(H^2)
\le\frac{Fr}{2}\sum A_\ell n_\ell^2+O(H^2).
\end{equation}
The assembled $v$ lies in $\operatorname{range}P$ and
$\sum_j(g_jf)^\top v_j=0$ exactly.

\paragraph{Observable dual and energy costs.}
Writing $T$ for the stacked whitening map, the real direction $Tv$
acts on the whitened conditional innovations and has noise variance
$\|v\|^2$. Fenchel's inequality gives
$G\ge a_f^\top v-\|v\|^2/2-h_{\mathcal Z}(\lambda)$, thus
\begin{equation}\label{eq:dual}
2D_{K,H}(\pi)\le
\|(I-P)a_f\|^2+\|Pa_f-v\|^2+2h_{\mathcal Z}(\lambda).
\end{equation}
On active inputs, auxiliary scalar errors are at most $(\eta+r)(n_\ell+k)$.
The at most four tail errors are $O(H)$, prefix errors zero, and inactive
blocks contribute $O(1)$.
Therefore $\|a_f-U\|^2=O(H^2)$ and, by contraction and repair,
\[
\|Pa_f-v\|^2\le2\|a_f-U\|^2+2C_k^2F\sum A_\ell^2/n_\ell=O(H^2).
\]
The two input-gap centers lie at offsets $1/2\pm(n_\ell+k)/2$.
With fixed $\beta=\beta_K(k)$ and fixed $L_1,L_2$ from \eqref{eq:affine},
their exact oracle norm loss is
\[
2\beta(A_\ell+\eta/2)^2+\frac{\beta\eta^2}{2}(n_\ell+k)^2
+4\eta(A_\ell+\eta/2)L_1+2\eta^2L_2 .
\]
Consequently $\|(I-P)a_f\|^2=2\beta\sum A_\ell^2+O(H^{3/2})$.
The non-reflection cross has been retained, not asserted zero.

Substituting in \eqref{eq:support}--\eqref{eq:dual} gives
$2D\le2\beta\sum A_\ell^2+Fr\sum A_\ell n_\ell^2+O(H^2)$.
Power summation yields
$\sum A_\ell^2=\eta^2H^{5/2}/(5\sqrt\xi)+O(H^2)$ and
$\sum A_\ell n_\ell^2=\eta\sqrt\xi H^{5/2}/5+O(H^2)$.
This proves \eqref{eq:upper}. Minimizing at
$\xi_*=2\beta\eta/(Fr)$ matches \eqref{eq:lower}, proving \eqref{eq:sharp}.
\end{proof}

\begin{thebibliography}{10}
\bibitem{controlled}
S. Nitinawarat, G. K. Atia, and V. V. Veeravalli,
``Controlled sensing for multihypothesis testing,''
\emph{IEEE Trans. Autom. Control}, vol. 58, no. 10, pp. 2451--2464, 2013,
doi: 10.1109/TAC.2013.2261188.
\bibitem{markov}
S. Nitinawarat and V. V. Veeravalli,
``Controlled sensing for sequential multihypothesis testing with controlled
Markovian observations and non-uniform control cost,''
arXiv:1310.1844v2, 2014. [Online]. Available: \url{https://arxiv.org/abs/1310.1844}
\bibitem{zonotope}
J. K. Scott, R. Findeisen, R. D. Braatz, and D. M. Raimondo,
``Input design for guaranteed fault diagnosis using zonotopes,''
\emph{Automatica}, vol. 50, pp. 1580--1589, 2014.
\bibitem{convex}
A. Goldenshluger, A. Juditsky, and A. Nemirovski,
``Hypothesis testing by convex optimization,'' arXiv:1311.6765v7, 2016.
[Online]. Available: \url{https://arxiv.org/abs/1311.6765}
\bibitem{allan}
R. Schieder and C. Kramer,
``Optimization of radio astronomical observations using Allan variance measurements,''
arXiv:astro-ph/0105071, 2001.
[Online]. Available: \url{https://arxiv.org/abs/astro-ph/0105071}

\bibitem{kim}
K.-K. K. Kim, D. M. Raimondo, and R. D. Braatz,
``Optimum input design for fault detection and diagnosis: Model-based prediction
and statistical distance measures,'' in \emph{Proc. European Control Conf.},
2013, pp. 1940--1945.
\bibitem{closedset}
D. M. Raimondo, G. R. Marseglia, R. D. Braatz, and J. K. Scott,
``Closed-loop input design for guaranteed fault diagnosis using set-valued
observers,'' \emph{Automatica}, vol. 74, pp. 107--117, 2016,
doi: 10.1016/j.automatica.2016.07.033.
\bibitem{feedback2012}
A. Esna Ashari, R. Nikoukhah, and S. L. Campbell,
``Effects of feedback on active fault detection,'' \emph{Automatica},
vol. 48, pp. 866--872, 2012, doi: 10.1016/j.automatica.2012.02.020.
\bibitem{robustfeedback2012}
A. Esna Ashari, R. Nikoukhah, and S. L. Campbell,
``Active robust fault detection in closed-loop systems: Quadratic optimization
approach,'' \emph{IEEE Trans. Autom. Control}, vol. 57, no. 10,
pp. 2532--2544, 2012, doi: 10.1109/TAC.2012.2188430.
\bibitem{switchback}
I. Bojinov, D. Simchi-Levi, and J. Zhao,
``Design and analysis of switchback experiments,'' arXiv:2009.00148v4, 2025.
[Online]. Available: \url{https://arxiv.org/abs/2009.00148}

\end{thebibliography}
\end{document}
